% ====================================================================
\chapter{Diseño del Sistema}
\label{cap:diseno_sistema}
% ====================================================================

\section{Arquitectura del sistema}

El diseño adopta una arquitectura híbrida desacoplada. La extracción de puntos articulares y la clasificación gestual se ejecutan en el equipo del usuario, mientras que el módulo de traducción semántica opera como un servicio remoto distribuido.

\subsection{Justificación de la partición y modelo de privacidad}
La ejecución local de la extracción y clasificación gestual garantiza que la señal de video cruda y las coordenadas articulares no abandonen el equipo del usuario. La demanda de memoria RAM y capacidad de cómputo de los modelos de lenguaje exige situar el componente de traducción semántica en un servicio remoto en la nube. Hacia dicho servicio remoto se transmite únicamente la secuencia textual de glosas (en el canal directo) o la oración en español (en el canal inverso), sin transferir información biométrica o de video.

\subsection{Flujo integral de datos entre componentes}
La interacción entre el navegador, el motor local y los servicios remotos se organiza según la estructura representada en la Figura~\ref{fig:diagrama_arquitectura}. El procesamiento local gestiona la captura y clasificación, mientras que la red transmite exclusivamente mensajes en formato JSON.

\begin{figure}[!ht]
\centering
\resizebox{1\textwidth}{!}{%
\begin{circuitikz}
\tikzstyle{every node}=[font=\fontsize{8.0pt}{10.4pt}\selectfont]

% Caja Cliente
\draw (1.125,14.25) rectangle (5.375,7.25);
\node[anchor=north, align=center, inner xsep=0.080cm, inner ysep=0.085cm, rounded corners=0.020cm, font=\fontsize{11.9pt}{15.5pt}\selectfont] at (3.25,14.1) {Cliente};

\draw (1.875,13.25) rectangle (4.375,12.5);
\node[anchor=north, align=center, inner xsep=0.080cm, inner ysep=0.085cm, rounded corners=0.020cm, font=\fontsize{10.8pt}{14.1pt}\selectfont] at (3.125,13.1) {Google Meet};

\draw (1.875,11.375) rectangle (4.375,10.625);
\node[anchor=north, align=center, inner xsep=0.080cm, inner ysep=0.085cm, rounded corners=0.020cm, font=\fontsize{10.8pt}{14.1pt}\selectfont] at (3.125,11.225) {Extension};

\draw (1.75,8.875) rectangle (4.5,7.75);
\node[anchor=center, align=center, inner xsep=0.080cm, inner ysep=0.085cm, rounded corners=0.020cm, font=\fontsize{10.8pt}{14.1pt}\selectfont] at (3.125,8.3125) {Motor ILSA. \\ Clasificador};

% Flechas Cliente
\draw [-{Stealth[scale=1.5]}] (3.25,12.5) -- (3.25,11.375) node[pos=0.45, fill=white, fill opacity=1, text opacity=1, inner xsep=0.080cm, inner ysep=0.085cm, rounded corners=0.020cm]{fotogramas};

\draw [-{Stealth[scale=1.5]}] (2.375,10.625) -- (2.375,8.875) node[pos=0.5, align=center, fill=white, fill opacity=1, text opacity=1, inner xsep=0.080cm, inner ysep=0.085cm, rounded corners=0.020cm]{frames \\ coordenadas \\ audio};

\draw [-{Stealth[scale=1.5]}] (3.75,8.875) -- (3.75,10.625) node[pos=0.5, fill=white, fill opacity=1, text opacity=1, inner xsep=0.080cm, inner ysep=0.085cm, rounded corners=0.020cm]{subtitulos};

% Caja Nube
\draw (7.875,11.5) rectangle (11.5,7.25);
\node[anchor=north, align=center, inner xsep=0.080cm, inner ysep=0.085cm, rounded corners=0.020cm, font=\fontsize{10.8pt}{14.1pt}\selectfont] at (9.6875,11.35) {Nube};

\draw (8.375,10.5) rectangle (11,9.25);
\node[anchor=north, align=center, inner xsep=0.080cm, inner ysep=0.085cm, rounded corners=0.020cm, font=\fontsize{10.8pt}{14.1pt}\selectfont] at (9.6875,10.35) {Servicio \\ semantico};

\draw (8.375,8.875) rectangle (11,7.625);
\node[anchor=north, align=center, inner xsep=0.080cm, inner ysep=0.085cm, rounded corners=0.020cm, font=\fontsize{10.8pt}{14.1pt}\selectfont] at (9.6875,8.725) {Servicio de \\ poses};

% Flechas Cliente <-> Nube
\draw [-{Stealth[scale=1.5]}] (4.5,8.25) -- (8.375,8.25) node[pos=0.5, fill=white, fill opacity=1, text opacity=1, inner xsep=0.080cm, inner ysep=0.085cm, rounded corners=0.020cm]{nombre de glosa};

\draw [-{Stealth[scale=1.5]}] (8.375,7.875) -- (4.5,7.875) node[pos=0.5, fill=white, fill opacity=1, text opacity=1, inner xsep=0.080cm, inner ysep=0.085cm, rounded corners=0.020cm]{archivo pose};

\draw [{Stealth[scale=1.5]}-{Stealth[scale=1.5]}] (4.5,8.625) -- (8.375,9.875) node[pos=0.5, fill=white, fill opacity=1, text opacity=1, inner xsep=0.080cm, inner ysep=0.085cm, rounded corners=0.020cm]{glosas / oración};

\end{circuitikz}%
}
\caption{Diagrama de la arquitectura general del sistema.}
\label{fig:diagrama_arquitectura}
\end{figure}

Para detallar la secuencia de intercambio de mensajes, las Figuras~\ref{fig:diagrama_secuencia_directo} y \ref{fig:diagrama_secuencia_inverso} ilustran el comportamiento de los canales directo e inverso respectivamente.

\begin{figure}[!ht]
\centering
\resizebox{1\textwidth}{!}{%
\begin{circuitikz}
\tikzstyle{every node}=[font=\fontsize{8.0pt}{10.4pt}\selectfont]

% Google Meet (Caja superior e inferior)
\draw (-2.25,16) rectangle (-0.625,15.125);
\node[anchor=center, align=center, inner xsep=0.080cm, inner ysep=0.085cm, rounded corners=0.020cm, font=\fontsize{8.0pt}{10.4pt}\selectfont] at (-1.4375,15.5625) {Google \\ Meet};
\draw [short] (-1.5,15.125) -- (-1.5,3.125);

\draw (-2.25,3.125) rectangle (-0.625,2.25);
\node[anchor=center, align=center, inner xsep=0.080cm, inner ysep=0.085cm, rounded corners=0.020cm, font=\fontsize{8.0pt}{10.4pt}\selectfont] at (-1.4375,2.6875) {Google \\ Meet};

% Extensión (Caja superior e inferior)
\draw (0.625,16) rectangle (2.25,15.125);
\node[anchor=center, align=center, inner xsep=0.080cm, inner ysep=0.085cm, rounded corners=0.020cm, font=\fontsize{8.0pt}{10.4pt}\selectfont] at (1.4375,15.5625) {Extensión};
\draw [short] (1.375,15.125) -- (1.375,3.125);

\draw (0.625,3.125) rectangle (2.25,2.25);
\node[anchor=center, align=center, inner xsep=0.080cm, inner ysep=0.085cm, rounded corners=0.020cm, font=\fontsize{8.0pt}{10.4pt}\selectfont] at (1.4375,2.6875) {Extensión};

% ILSA (Caja superior e inferior)
\draw (3.5,16) rectangle (5.125,15.125);
\node[anchor=center, align=center, inner xsep=0.080cm, inner ysep=0.085cm, rounded corners=0.020cm, font=\fontsize{8.0pt}{10.4pt}\selectfont] at (4.3125,15.5625) {ILSA};
\draw [short] (4.25,15.125) -- (4.25,3.125) node[pos=0.49, fill=white, fill opacity=1, text opacity=1, inner xsep=0.080cm, inner ysep=0.085cm, rounded corners=0.020cm]{cierre por temporizador};

\draw (3.5,3.125) rectangle (5.125,2.25);
\node[anchor=center, align=center, inner xsep=0.080cm, inner ysep=0.085cm, rounded corners=0.020cm, font=\fontsize{8.0pt}{10.4pt}\selectfont] at (4.3125,2.6875) {ILSA};

% Servidor semántico (Caja superior e inferior)
\draw (6.375,16) rectangle (8,15.125);
\node[anchor=center, align=center, inner xsep=0.080cm, inner ysep=0.085cm, rounded corners=0.020cm, font=\fontsize{8.0pt}{10.4pt}\selectfont] at (7.1875,15.5625) {Servidor \\ semántico};
\draw [short] (7.125,15.125) -- (7.125,3.125);

\draw (6.375,3.125) rectangle (8,2.25);
\node[anchor=center, align=center, inner xsep=0.080cm, inner ysep=0.085cm, rounded corners=0.020cm, font=\fontsize{8.0pt}{10.4pt}\selectfont] at (7.1875,2.6875) {Servidor \\ semántico};

% Loop por seña
\draw [dashed] (-1.875,15) rectangle (5.75,10.5);
\node[anchor=north east, align=right, inner xsep=0.080cm, inner ysep=0.085cm, rounded corners=0.020cm, font=\fontsize{8.0pt}{10.4pt}\selectfont] at (5.6,14.85) {[loop por seña]};

% Flujos de interacción
\draw [-{Stealth[scale=1.5]}] (-1.5,14) -- (1.375,14) node[pos=0.5, above, fill=white, fill opacity=1, text opacity=1, inner xsep=0.080cm, inner ysep=0.085cm, rounded corners=0.020cm]{video interceptado};

\draw [draw opacity=0, fill={rgb,255:red,215; green,215; blue,215}, fill opacity=1] (0.25,13.75) rectangle (3,13.25);
\node[anchor=center, align=center, inner xsep=0.080cm, inner ysep=0.085cm, rounded corners=0.020cm, font=\fontsize{8.0pt}{10.4pt}\selectfont] at (1.625,13.5) {extrae pose y manos};

\draw [-{Stealth[scale=1.5]}] (1.375,12.5) -- (4.25,12.5) node[pos=0.5, above, fill=white, fill opacity=1, text opacity=1, inner xsep=0.080cm, inner ysep=0.085cm, rounded corners=0.020cm]{POST /sign};

\draw [draw opacity=0, fill={rgb,255:red,215; green,215; blue,215}, fill opacity=1] (2.875,12.25) rectangle (5.625,11.75);
\node[anchor=center, align=center, inner xsep=0.080cm, inner ysep=0.085cm, rounded corners=0.020cm, font=\fontsize{8.0pt}{10.4pt}\selectfont] at (4.25,12) {clasificador y filtro};

\draw [-{Stealth[scale=1.5]}, dash pattern=on 6pt off 2pt] (4.25,11.375) -- (1.375,11.375) node[pos=0.5, fill=white, fill opacity=1, text opacity=1, inner xsep=0.080cm, inner ysep=0.085cm, rounded corners=0.020cm]{200 OK};

\draw [draw opacity=0, fill={rgb,255:red,215; green,215; blue,215}, fill opacity=1] (2.875,11.125) rectangle (5.625,10.625);
\node[anchor=center, align=center, inner xsep=0.080cm, inner ysep=0.085cm, rounded corners=0.020cm, font=\fontsize{8.0pt}{10.4pt}\selectfont] at (4.25,10.875) {acumula glosa};

\draw [draw opacity=0, fill={rgb,255:red,215; green,215; blue,215}, fill opacity=1] (-4.625,10.25) rectangle (4.625,9.75);
\node[anchor=center, align=center, inner xsep=0.080cm, inner ysep=0.085cm, rounded corners=0.020cm, font=\fontsize{8.0pt}{10.4pt}\selectfont] at (0,10) {pausa física de inactividad};

\draw [-{Stealth[scale=1.5]}] (4.25,8.25) -- (7.125,8.25) node[pos=0.5, above, fill=white, fill opacity=1, text opacity=1, inner xsep=0.080cm, inner ysep=0.085cm, rounded corners=0.020cm]{POST /sign};

\draw [draw opacity=0, fill={rgb,255:red,215; green,215; blue,215}, fill opacity=1] (5.75,8) rectangle (8.5,7.5);
\node[anchor=center, align=center, inner xsep=0.080cm, inner ysep=0.085cm, rounded corners=0.020cm, font=\fontsize{8.0pt}{10.4pt}\selectfont] at (7.125,7.75) {interpretación};

\draw [-{Stealth[scale=1.5]}, dash pattern=on 6pt off 2pt] (7.125,6.625) -- (4.25,6.625) node[pos=0.5, above, fill=white, fill opacity=1, text opacity=1, inner xsep=0.080cm, inner ysep=0.085cm, rounded corners=0.020cm]{200 OK \{"text": "..."\}};

\draw [-{Stealth[scale=1.5]}] (4.25,6.25) -- (1.375,6.25) node[pos=0.5, above, fill=white, fill opacity=1, text opacity=1, inner xsep=0.080cm, inner ysep=0.085cm, rounded corners=0.020cm]{oración};

\draw [-{Stealth[scale=1.5]}] (1.375,5.875) -- (-1.5,5.875) node[pos=0.5, above, fill=white, fill opacity=1, text opacity=1, inner xsep=0.080cm, inner ysep=0.085cm, rounded corners=0.020cm]{oración};

\draw [-{Stealth[scale=1.5]}] (-1.5,5.125) -- (10.125,5.125) node[pos=0.5, above, fill=white, fill opacity=1, text opacity=1, inner xsep=0.080cm, inner ysep=0.085cm, rounded corners=0.020cm]{Muestras subtítulos};

\draw [draw opacity=0, fill={rgb,255:red,215; green,215; blue,215}, fill opacity=1] (2.875,4.75) rectangle (5.625,4.25);
\node[anchor=center, align=center, inner xsep=0.080cm, inner ysep=0.085cm, rounded corners=0.020cm, font=\fontsize{8.0pt}{10.4pt}\selectfont] at (4.25,4.5) {síntesis de voz};

\draw [-{Stealth[scale=1.5]}] (4.25,3.75) -- (10.125,3.75) node[pos=0.5, above, fill=white, fill opacity=1, text opacity=1, inner xsep=0.080cm, inner ysep=0.085cm, rounded corners=0.020cm]{voz};

\draw [-{Stealth[scale=1.5]}] (-4.375,14.125) -- (-1.5,14.125) node[pos=0.5, above, fill=white, fill opacity=1, text opacity=1, inner xsep=0.080cm, inner ysep=0.085cm, rounded corners=0.020cm]{movimiento};

% Usuario señante (Izquierda)
\draw [short] (-4.375,15.125) -- (-4.375,3.625);
\draw (-4.375,16.125) circle (0.125cm);
\draw [short] (-4.375,16) -- (-4.375,15.75);
\draw [short] (-4.375,15.75) -- (-4.5,15.625);
\draw [short] (-4.375,15.75) -- (-4.25,15.625);
\draw [short] (-4.5,15.875) -- (-4.25,15.875);
\node[anchor=center, align=center, inner xsep=0.080cm, inner ysep=0.085cm, rounded corners=0.020cm, font=\fontsize{6.8pt}{8.9pt}\selectfont] at (-4.3125,15.3125) {Usuario señante};

\draw (-4.375,3.25) circle (0.125cm);
\draw [short] (-4.375,3.125) -- (-4.375,2.875);
\draw [short] (-4.375,2.875) -- (-4.5,2.75);
\draw [short] (-4.375,2.875) -- (-4.25,2.75);
\draw [short] (-4.5,3) -- (-4.25,3);
\node[anchor=center, align=center, inner xsep=0.080cm, inner ysep=0.085cm, rounded corners=0.020cm, font=\fontsize{6.8pt}{8.9pt}\selectfont] at (-4.3125,2.4375) {Usuario señante};

% Usuario oyente (Derecha)
\draw [short] (10.125,15.125) -- (10.125,3.5);
\draw (10.125,16.125) circle (0.125cm);
\draw [short] (10.125,16) -- (10.125,15.75);
\draw [short] (10.125,15.75) -- (10,15.625);
\draw [short] (10.125,15.75) -- (10.25,15.625);
\draw [short] (10,15.875) -- (10.25,15.875);
\node[anchor=center, align=center, inner xsep=0.080cm, inner ysep=0.085cm, rounded corners=0.020cm, font=\fontsize{6.8pt}{8.9pt}\selectfont] at (10.1875,15.3125) {Usuario oyente};

\draw (10.125,3.125) circle (0.125cm);
\draw [short] (10.125,3) -- (10.125,2.75);
\draw [short] (10.125,2.75) -- (10,2.625);
\draw [short] (10.125,2.75) -- (10.25,2.625);
\draw [short] (10,2.875) -- (10.25,2.875);
\node[anchor=center, align=center, inner xsep=0.080cm, inner ysep=0.085cm, rounded corners=0.020cm, font=\fontsize{6.8pt}{8.9pt}\selectfont] at (10.1875,2.3125) {Usuario oyente};

\end{circuitikz}%
}
\caption{Diagrama de secuencia del flujo de traducción de Lengua de Señas a voz y subtítulos.}
\label{fig:diagrama_secuencia_directo}
\end{figure}

\begin{figure}[!ht]
\centering
\resizebox{1\textwidth}{!}{%
\begin{circuitikz}
\tikzstyle{every node}=[font=\fontsize{8.0pt}{10.4pt}\selectfont]

% Google Meet (Cajas superior e inferior)
\draw (-1.125,16.75) rectangle (0.5,15.875);
\node[anchor=center, align=center, inner xsep=0.080cm, inner ysep=0.085cm, rounded corners=0.020cm, font=\fontsize{8.0pt}{10.4pt}\selectfont] at (-0.3125,16.3125) {Google \\ Meet};
\draw [short] (-0.375,15.875) -- (-0.375,2.625);

\draw (-1.125,2.625) rectangle (0.5,1.75);
\node[anchor=center, align=center, inner xsep=0.080cm, inner ysep=0.085cm, rounded corners=0.020cm, font=\fontsize{8.0pt}{10.4pt}\selectfont] at (-0.3125,2.1875) {Google \\ Meet};

% Extensión (Cajas superior e inferior)
\draw (1.75,16.75) rectangle (3.375,15.875);
\node[anchor=center, align=center, inner xsep=0.080cm, inner ysep=0.085cm, rounded corners=0.020cm, font=\fontsize{8.0pt}{10.4pt}\selectfont] at (2.5625,16.3125) {Extensión};
\draw [short] (2.5,15.875) -- (2.5,2.625);

\draw (1.75,2.625) rectangle (3.375,1.75);
\node[anchor=center, align=center, inner xsep=0.080cm, inner ysep=0.085cm, rounded corners=0.020cm, font=\fontsize{8.0pt}{10.4pt}\selectfont] at (2.5625,2.1875) {Extensión};

% ILSA (Cajas superior e inferior)
\draw (4.625,16.75) rectangle (6.25,15.875);
\node[anchor=center, align=center, inner xsep=0.080cm, inner ysep=0.085cm, rounded corners=0.020cm, font=\fontsize{8.0pt}{10.4pt}\selectfont] at (5.4375,16.3125) {ILSA};
\draw [short] (5.375,15.875) -- (5.375,2.625);

\draw (4.625,2.625) rectangle (6.25,1.75);
\node[anchor=center, align=center, inner xsep=0.080cm, inner ysep=0.085cm, rounded corners=0.020cm, font=\fontsize{8.0pt}{10.4pt}\selectfont] at (5.4375,2.1875) {ILSA};

% Servidor semántico (Cajas superior e inferior)
\draw (7.5,16.75) rectangle (9.125,15.875);
\node[anchor=center, align=center, inner xsep=0.080cm, inner ysep=0.085cm, rounded corners=0.020cm, font=\fontsize{8.0pt}{10.4pt}\selectfont] at (8.3125,16.3125) {Servidor \\ semántico};
\draw [short] (8.25,15.875) -- (8.25,2.625);

\draw (7.5,2.625) rectangle (9.125,1.75);
\node[anchor=center, align=center, inner xsep=0.080cm, inner ysep=0.085cm, rounded corners=0.020cm, font=\fontsize{8.0pt}{10.4pt}\selectfont] at (8.3125,2.1875) {Servidor \\ semántico};

% Servidor de poses (Cajas superior e inferior)
\draw (10.5,16.75) rectangle (12.125,15.875);
\node[anchor=center, align=center, inner xsep=0.080cm, inner ysep=0.085cm, rounded corners=0.020cm, font=\fontsize{8.0pt}{10.4pt}\selectfont] at (11.3125,16.3125) {Servidor \\ de poses};
\draw [short] (11.25,15.875) -- (11.25,2.625);

\draw (10.5,2.625) rectangle (12.125,1.75);
\node[anchor=center, align=center, inner xsep=0.080cm, inner ysep=0.085cm, rounded corners=0.020cm, font=\fontsize{8.0pt}{10.4pt}\selectfont] at (11.3125,2.1875) {Servidor \\ de poses};

% Usuario oyente (Izquierda)
\draw [short] (-3.25,15.875) -- (-3.25,3);
\draw (-3.25,16.875) circle (0.125cm);
\draw [short] (-3.25,16.75) -- (-3.25,16.5);
\draw [short] (-3.25,16.5) -- (-3.375,16.375);
\draw [short] (-3.25,16.5) -- (-3.125,16.375);
\draw [short] (-3.375,16.625) -- (-3.125,16.625);
\node[anchor=center, align=center, inner xsep=0.080cm, inner ysep=0.085cm, rounded corners=0.020cm, font=\fontsize{6.8pt}{8.9pt}\selectfont] at (-3.1875,16.0625) {Usuario oyente};

\draw (-3.25,2.75) circle (0.125cm);
\draw [short] (-3.25,2.625) -- (-3.25,2.375);
\draw [short] (-3.25,2.375) -- (-3.375,2.25);
\draw [short] (-3.25,2.375) -- (-3.125,2.25);
\draw [short] (-3.375,2.5) -- (-3.125,2.5);
\node[anchor=center, align=center, inner xsep=0.080cm, inner ysep=0.085cm, rounded corners=0.020cm, font=\fontsize{6.8pt}{8.9pt}\selectfont] at (-3.1875,1.9375) {Usuario oyente};

% Usuario señante (Derecha)
\draw [short] (14.25,15.875) -- (14.25,2.875);
\draw (14.25,16.875) circle (0.125cm);
\draw [short] (14.25,16.75) -- (14.25,16.5);
\draw [short] (14.25,16.5) -- (14.125,16.375);
\draw [short] (14.25,16.5) -- (14.375,16.375);
\draw [short] (14.125,16.625) -- (14.375,16.625);
\node[anchor=center, align=center, inner xsep=0.080cm, inner ysep=0.085cm, rounded corners=0.020cm, font=\fontsize{6.8pt}{8.9pt}\selectfont] at (14.3125,16.0625) {Usuario señante};

\draw (14.25,2.625) circle (0.125cm);
\draw [short] (14.25,2.5) -- (14.25,2.25);
\draw [short] (14.25,2.25) -- (14.125,2.125);
\draw [short] (14.25,2.25) -- (14.375,2.125);
\draw [short] (14.125,2.375) -- (14.375,2.375);
\node[anchor=center, align=center, inner xsep=0.080cm, inner ysep=0.085cm, rounded corners=0.020cm, font=\fontsize{6.8pt}{8.9pt}\selectfont] at (14.3125,1.8125) {Usuario señante};

% Marcos de agrupación condicionales y de bucles
\draw [dash pattern=on 6pt off 2pt] (2.125,12.125) rectangle (5.75,14.5);
\node[anchor=north east, align=right, inner xsep=0.080cm, inner ysep=0.085cm, rounded corners=0.020cm, font=\fontsize{8.0pt}{10.4pt}\selectfont] at (5.6,14.35) {[Si viene de micrófono]};

\draw [dash pattern=on 6pt off 2pt] (4.875,4.375) rectangle (11.875,9.25);
\node[anchor=north east, align=right, inner xsep=0.080cm, inner ysep=0.085cm, rounded corners=0.020cm, font=\fontsize{8.0pt}{10.4pt}\selectfont] at (11.725,9.1) {[Por cada glosa]};

\draw [dash pattern=on 6pt off 2pt] (5.125,5.625) rectangle (11.5,8);
\node[anchor=north east, align=right, inner xsep=0.080cm, inner ysep=0.085cm, rounded corners=0.020cm, font=\fontsize{8.0pt}{10.4pt}\selectfont] at (11.35,7.85) {[No existe archivo]};

% Cuadro de proceso de interpretación
\draw [draw opacity=0, fill={rgb,255:red,215; green,215; blue,215}, fill opacity=1] (6.75,10.375) rectangle (9.5,10.875);
\node[anchor=center, align=center, inner xsep=0.080cm, inner ysep=0.085cm, rounded corners=0.020cm, font=\fontsize{8.0pt}{10.4pt}\selectfont] at (8.125,10.625) {interpretación};

% Flujos de mensajes (Flechas)
\draw [-{Stealth[scale=1.5]}] (-3.25,14.875) -- (-0.375,14.875) node[pos=0.5, above, fill=white, fill opacity=1, text opacity=1, inner xsep=0.080cm, inner ysep=0.085cm, rounded corners=0.020cm]{habla o escribe};

\draw [-{Stealth[scale=1.5]}] (-0.375,14.75) -- (2.5,14.75) node[pos=0.5, above, fill=white, fill opacity=1, text opacity=1, inner xsep=0.080cm, inner ysep=0.085cm, rounded corners=0.020cm]{audio/texto chat};

\draw [-{Stealth[scale=1.5]}] (2.5,13.25) -- (5.375,13.25) node[pos=0.5, above, fill=white, fill opacity=1, text opacity=1, inner xsep=0.080cm, inner ysep=0.085cm, rounded corners=0.020cm]{POST /stt};

\draw [-{Stealth[scale=1.5]}, dash pattern=on 6pt off 2pt] (5.375,12.5) -- (2.5,12.5) node[pos=0.5, above, fill=white, fill opacity=1, text opacity=1, inner xsep=0.080cm, inner ysep=0.085cm, rounded corners=0.020cm]{oración};

\draw [-{Stealth[scale=1.5]}] (2.5,11.375) -- (5.375,11.375) node[pos=0.5, above, fill=white, fill opacity=1, text opacity=1, inner xsep=0.080cm, inner ysep=0.085cm, rounded corners=0.020cm]{POST /hearing};

\draw [-{Stealth[scale=1.5]}] (5.375,11.125) -- (8.25,11.125) node[pos=0.5, above, fill=white, fill opacity=1, text opacity=1, inner xsep=0.080cm, inner ysep=0.085cm, rounded corners=0.020cm]{POST /hearing};

\draw [-{Stealth[scale=1.5]}, dash pattern=on 6pt off 2pt] (8.25,9.5) -- (5.375,9.5) node[pos=0.5, above, fill=white, fill opacity=1, text opacity=1, inner xsep=0.080cm, inner ysep=0.085cm, rounded corners=0.020cm]{200 OK, glosas};

\draw [-{Stealth[scale=1.5]}] (5.375,8.375) -- (11.25,8.375) node[pos=0.5, above, fill=white, fill opacity=1, text opacity=1, inner xsep=0.080cm, inner ysep=0.085cm, rounded corners=0.020cm]{GET /tposes/\{archivo\}};

\draw [-{Stealth[scale=1.5]}, dash pattern=on 6pt off 2pt] (11.25,6.75) -- (5.375,6.75) node[pos=0.5, above, fill=white, fill opacity=1, text opacity=1, inner xsep=0.080cm, inner ysep=0.085cm, rounded corners=0.020cm]{404 NOT FOUND};

\draw [-{Stealth[scale=1.5]}] (5.375,6) -- (11.25,6) node[pos=0.5, above, fill=white, fill opacity=1, text opacity=1, inner xsep=0.080cm, inner ysep=0.085cm, rounded corners=0.020cm]{GET /fingerspell?letters=...};

\draw [-{Stealth[scale=1.5]}, dash pattern=on 6pt off 2pt] (11.25,4.75) -- (5.375,4.75) node[pos=0.5, above, fill=white, fill opacity=1, text opacity=1, inner xsep=0.080cm, inner ysep=0.085cm, rounded corners=0.020cm]{200, archivo/s de seña/s};

\draw [-{Stealth[scale=1.5]}, dash pattern=on 6pt off 2pt] (5.375,4) -- (2.5,4) node[pos=0.5, above, fill=white, fill opacity=1, text opacity=1, inner xsep=0.080cm, inner ysep=0.085cm, rounded corners=0.020cm]{Glosas y animación};

\draw [-{Stealth[scale=1.5]}] (2.5,3.75) -- (-0.375,3.75) node[pos=0.5, above, fill=white, fill opacity=1, text opacity=1, inner xsep=0.080cm, inner ysep=0.085cm, rounded corners=0.020cm]{Subtítulo/esqueleto};

\draw [-{Stealth[scale=1.5]}] (-0.375,3.375) -- (14.25,3.375) node[pos=0.5, above, fill=white, fill opacity=1, text opacity=1, inner xsep=0.080cm, inner ysep=0.085cm, rounded corners=0.020cm]{Avatar señando / subtítulos};

\end{circuitikz}%
}
\caption{Diagrama de secuencia del flujo de interpretación de voz/texto a señas para el usuario señante.}
\label{fig:diagrama_secuencia_inverso}
\end{figure}

\section{Diseño de contratos de API REST}

La comunicación entre el cliente local y los microservicios se rige por contratos JSON estructurados.

\subsection{API local del motor ILSA}
El motor ILSA expone una interfaz HTTP en el bucle local (\texttt{127.0.0.1:8765}). La extensión web actúa como cliente y el proceso local como servidor. El esquema general de componentes internos se representa en la Figura~\ref{fig:ilsa_local_arquitectura}.

\begin{figure}[!ht]
\centering
\resizebox{1\textwidth}{!}{%
\begin{circuitikz}
\tikzstyle{every node}=[font=\fontsize{8.0pt}{10.4pt}\selectfont]

% Contenedor principal: ILSA (motor local)
\draw [line width=0.6pt] (0.625,15) rectangle (12.625,4.875);
\node[anchor=north west, align=left, inner xsep=0.080cm, inner ysep=0.085cm, rounded corners=0.020cm, font=\fontsize{10.8pt}{14.1pt}\selectfont] at (0.775,14.85) {ILSA (motor local)};

% Caja Extension
\draw (0.5,16.875) rectangle (12.625,16.125);
\node[anchor=north, align=center, inner xsep=0.080cm, inner ysep=0.085cm, rounded corners=0.020cm, font=\fontsize{10.8pt}{14.1pt}\selectfont] at (6.5625,16.725) {Extension};

% Subcaja Clasificador
\draw [line width=0.7pt] (1,11) rectangle (4.625,5.25);
\node[anchor=north west, align=left, inner xsep=0.080cm, inner ysep=0.085cm, rounded corners=0.020cm, font=\fontsize{10.8pt}{14.1pt}\selectfont] at (1.15,10.85) {clasificador};

\draw (1.25,10) rectangle (4.375,8.75);
\node[anchor=north, align=center, inner xsep=0.080cm, inner ysep=0.085cm, rounded corners=0.020cm, font=\fontsize{10.8pt}{14.1pt}\selectfont] at (2.8125,9.85) {extraccion de \\ puntos corporales};

\draw (1.25,7.875) rectangle (4.375,7.125);
\node[anchor=north, align=center, inner xsep=0.080cm, inner ysep=0.085cm, rounded corners=0.020cm, font=\fontsize{10.8pt}{14.1pt}\selectfont] at (2.8125,7.725) {clasificador};

\draw (1.25,6.25) rectangle (4.375,5.5);
\node[anchor=north, align=center, inner xsep=0.080cm, inner ysep=0.085cm, rounded corners=0.020cm, font=\fontsize{10.8pt}{14.1pt}\selectfont] at (2.8125,6.1) {filtros y buffer};

% Cajas de síntesis y sesión
\draw (9.875,11) rectangle (12.375,9.875);
\node[anchor=center, align=center, inner xsep=0.080cm, inner ysep=0.085cm, rounded corners=0.020cm, font=\fontsize{10.8pt}{14.1pt}\selectfont] at (11.125,10.4375) {sintesis de \\ voz};

\draw (5.125,11) rectangle (9.375,9.875);
\node[anchor=center, align=center, inner xsep=0.080cm, inner ysep=0.085cm, rounded corners=0.020cm, font=\fontsize{10.8pt}{14.1pt}\selectfont] at (7.25,10.4375) {sesion oyente y avatar};

% Conexiones internas
\draw [-{Stealth[scale=1.5]}] (6,16.125) -- (6,11) node[pos=0.63, fill=white, fill opacity=1, text opacity=1, inner xsep=0.080cm, inner ysep=0.085cm, rounded corners=0.020cm]{POST /hearing};
\draw [short] (6.5,11.875) -- (6.5,11.875);
\draw [short] (4.25,12.25) -- (4.25,12.25);
\draw [short] (3.5,11.75) -- (3.5,11.75);

\draw [-{Stealth[scale=1.5]}] (2.75,8.75) -- (2.75,7.875) node[pos=0.41, fill=white, fill opacity=1, text opacity=1, inner xsep=0.080cm, inner ysep=0.085cm, rounded corners=0.020cm]{puntos corporales};
\draw [-{Stealth[scale=1.5]}] (2.75,7.125) -- (2.75,6.25) node[pos=0.38, fill=white, fill opacity=1, text opacity=1, inner xsep=0.080cm, inner ysep=0.085cm, rounded corners=0.020cm]{glosa y confianza};

% Contenedor Nube
\draw (0.625,3.75) rectangle (11.5,3);
\node[anchor=north, align=center, inner xsep=0.080cm, inner ysep=0.085cm, rounded corners=0.020cm, font=\fontsize{10.8pt}{14.1pt}\selectfont] at (6.0625,3.6) {Nube};

% Flujos con la Nube y la Extensión
\draw [-{Stealth[scale=1.5]}] (5.5,9.875) -- (5.5,3.75) node[pos=0.92, fill=white, fill opacity=1, text opacity=1, inner xsep=0.080cm, inner ysep=0.085cm, rounded corners=0.020cm]{GET /tposes/\{archivo\}};
\draw [-{Stealth[scale=1.5]}, dash pattern=on 6pt off 2pt] (6.75,3.75) -- (6.75,9.875) node[pos=0.6, fill=white, fill opacity=1, text opacity=1, inner xsep=0.080cm, inner ysep=0.085cm, rounded corners=0.020cm]{archivo pose};

\draw [-{Stealth[scale=1.5]}] (5.125,16.125) -- (5.125,13.625) -- (3.25,13.625) -- (3.25,10);
\node[fill=white, fill opacity=1, text opacity=1, inner xsep=0.080cm, inner ysep=0.085cm, rounded corners=0.020cm, font=\fontsize{8.0pt}{10.4pt}\selectfont] at (4.125,13.625) {POST /sign};

\draw [-{Stealth[scale=1.5]}] (11,3.75) -- (11,9.875) node[pos=0.5, fill=white, fill opacity=1, text opacity=1, inner xsep=0.080cm, inner ysep=0.085cm, rounded corners=0.020cm]{oración};
\draw [-{Stealth[scale=1.5]}, dash pattern=on 6pt off 2pt] (6.75,11) -- (6.75,16.125) node[pos=0.61, fill=white, fill opacity=1, text opacity=1, inner xsep=0.080cm, inner ysep=0.085cm, rounded corners=0.020cm]{animación};

\draw [-{Stealth[scale=1.5]}] (1.5,5.5) -- (1.5,3.75) node[pos=0.64, fill=white, fill opacity=1, text opacity=1, inner xsep=0.080cm, inner ysep=0.085cm, rounded corners=0.020cm]{POST /sign};
\draw [-{Stealth[scale=1.5]}, dash pattern=on 6pt off 2pt] (3.125,3.75) -- (3.125,5.5) node[pos=0.37, fill=white, fill opacity=1, text opacity=1, inner xsep=0.080cm, inner ysep=0.085cm, rounded corners=0.020cm]{oración};

\draw [dash pattern=on 6pt off 2pt, -{Stealth[scale=1.5]}] (1.25,5.75) -- (-1.375,5.75) -- (-1.375,16.5) -- (0.5,16.5);
\node[fill=white, fill opacity=1, text opacity=1, inner xsep=0.080cm, inner ysep=0.085cm, rounded corners=0.020cm, font=\fontsize{8.0pt}{10.4pt}\selectfont] at (-1.375,8.75) {oración};

\draw [-{Stealth[scale=1.5]}] (8,9.875) -- (8,3.75) node[pos=0.59, fill=white, fill opacity=1, text opacity=1, inner xsep=0.080cm, inner ysep=0.085cm, rounded corners=0.020cm]{POST /hearing};
\draw [-{Stealth[scale=1.5]}, dash pattern=on 6pt off 2pt] (9,3.75) -- (9,9.875) node[pos=0.53, fill=white, fill opacity=1, text opacity=1, inner xsep=0.080cm, inner ysep=0.085cm, rounded corners=0.020cm]{glosas};

\draw [dash pattern=on 6pt off 2pt, -{Stealth[scale=1.5]}] (11,11) -- (11,14.25) -- (11,16.125);
\node[fill=white, fill opacity=1, text opacity=1, inner xsep=0.080cm, inner ysep=0.085cm, rounded corners=0.020cm, font=\fontsize{8.0pt}{10.4pt}\selectfont] at (11,13.625) {audio};

\draw [-{Stealth[scale=1.5]}] (0.375,16.25) -- (-0.125,16.25) -- (-0.125,6) -- (1.25,6);
\node[fill=white, fill opacity=1, text opacity=1, inner xsep=0.080cm, inner ysep=0.085cm, rounded corners=0.020cm, align=center, font=\fontsize{8.0pt}{10.4pt}\selectfont] at (-0.125,12.625) {POST /session \\ POST /activity \\POST /utterance/end \\GET /health \\GET /state};

\draw [-{Stealth[scale=1.5]}] (8,16.125) -- (8,13.875) -- (8,12.375) -- (8,11);
\node[fill=white, fill opacity=1, text opacity=1, inner xsep=0.080cm, inner ysep=0.085cm, rounded corners=0.020cm, font=\fontsize{8.0pt}{10.4pt}\selectfont] at (8,12.375) {POST /stt};

\draw [dash pattern=on 6pt off 2pt, -{Stealth[scale=1.5]}] (8.875,11) -- (8.875,14.5) -- (8.875,16.125);
\node[fill=white, fill opacity=1, text opacity=1, inner xsep=0.080cm, inner ysep=0.085cm, rounded corners=0.020cm, font=\fontsize{8.0pt}{10.4pt}\selectfont] at (8.875,11.75) {oración};

\end{circuitikz}%
}
\caption{Arquitectura detallada del motor local ILSA y su interacción con los servicios en la nube.}
\label{fig:ilsa_local_arquitectura}
\end{figure}

El Cuadro~\ref{tab:api-ilsa-local} resume las rutas expuestas por el motor local.

\begin{table}[htbp]
  \centering
  \caption{Endpoints de la API local del motor ILSA (\texttt{http://127.0.0.1:8765}).}
  \label{tab:api-ilsa-local}
  \scriptsize
  \begin{tabular}{|p{0.10\textwidth}|p{0.24\textwidth}|p{0.48\textwidth}|p{0.12\textwidth}|}
    \hline
    \textbf{Método} & \textbf{Ruta} & \textbf{Contrato} & \textbf{Códigos} \\ \hline
    GET & \texttt{/health} & Disponibilidad del proceso, del modo, del clasificador y del traductor remoto. & 200 \\ \hline
    GET & \texttt{/config} & Umbrales de captura del modo en curso. & 200, 503 \\ \hline
    GET & \texttt{/state} & Búfer de glosas, último texto y estado del traductor. & 200, 503 \\ \hline
    GET & \texttt{/metrics} & Conteos y tiempos de la sesión en curso. & 200 \\ \hline
    POST & \texttt{/session} & Apertura de sesión. El cuerpo fija la lateralidad mediante \texttt{left\_handed}. & 200, 503 \\ \hline
    POST & \texttt{/sign} & Ingesta de hasta 60 fotogramas para clasificación en modo señante. & 200, 400, 503 \\ \hline
    POST & \texttt{/stt} & Transcripción de audio WAV en Base64 a texto en modo oyente. & 200, 400, 502 \\ \hline
    POST & \texttt{/hearing} & Envío de oración y recepción de animación en modo oyente. & 200, 400, 503 \\ \hline
    POST & \texttt{/activity} & Reinicio del temporizador de pausa de inactividad ($\sim 4$\,s). & 200, 503 \\ \hline
    POST & \texttt{/utterance/end} & Cierre explícito del enunciado y disparo de traducción. & 200, 503 \\ \hline
    GET & \texttt{/mediapipe/\{name\}} & Despacho de recursos estáticos del estimador visual. & 200, 400, 404 \\ \hline
    GET & \texttt{/download/exe} & Entrega del ejecutable del motor local. & 200, 404 \\ \hline
  \end{tabular}
\end{table}

\subsection{API del servicio semántico remoto}
\label{sec:api-semantico-remoto}
El microservicio remoto en la nube expone contratos REST sobre HTTPS para resolver la traducción semántica desacoplada. El Cuadro~\ref{tab:api-semantico-endpoints} resume los endpoints principales.

\begin{table}[htbp]
  \centering
  \caption{Endpoints del servicio semántico remoto.}
  \label{tab:api-semantico-endpoints}
  \begin{tabular}{@{}llp{8.2cm}@{}}
    \toprule
    Método & Ruta & Función \\
    \midrule
    GET & \texttt{/health} & Verificación de disponibilidad de los modelos de lenguaje en memoria. \\
    POST & \texttt{/sordo}, \texttt{/translate} & Canal directo: traducción de glosas LSA a oración en español rioplatense. \\
    POST & \texttt{/hearing} & Canal inverso: traducción de oración en español a glosas en estructura \textit{Topic-Comment}. \\
    \bottomrule
  \end{tabular}
\end{table}

\subsection{API del servicio de archivos del avatar}
\label{sec:api-avatar}
El subsistema de representación gestual de salida consulta un repositorio de archivos cinemáticos binarios. El Cuadro~\ref{tab:api-avatar-endpoints} detalla las rutas disponibles.

\begin{table}[htbp]
  \centering
  \caption{Endpoints del servicio de archivos del avatar.}
  \label{tab:api-avatar-endpoints}
  \begin{tabular}{@{}llp{7.6cm}@{}}
    \toprule
    Método & Ruta & Función \\
    \midrule
    GET & \texttt{/health} & Disponibilidad del proceso de archivos. \\
    GET & \texttt{/tposes/\{glosa\}.pose} & Entrega del flujo binario cinemático correspondiente a la glosa solicitada. \\
    GET & \texttt{/fingerspell} & Solicitud de clips del alfabeto manual para deletreo dactilológico. \\
    \bottomrule
  \end{tabular}
\end{table}

\section{Diseño de la representación gestual inversa}
\label{sec:repr_inversa}

El canal inverso procesa la lista de glosas devuelta por el traductor e infiere la secuencia cinemática correspondiente. Cada glosa se asocia a un clip de referencias cinemáticas previamente parametrizado.

\subsection{Procesamiento geométrico de coordenadas y segmentación temporal}
Los clips cinemáticos registran la secuencia de puntos clave extraídos mediante la ejecución previa e independiente de \textit{MediaPipe Holistic} sobre el video de un informante calificado. Las coordenadas planares $(x, y)$ se representan en el espacio discreto de píxeles de la cámara de captura ($x \in [0, W]$, $y \in [0, H]$), preservando las proporciones morfométricas nativas. La componente de profundidad $z$ posee carácter relativo con origen en la muñeca, actuando como métrica de jerarquía para resolver oclusiones mutuas. Por consiguiente, la proyección gráfica reproyecta $x$ e $y$ directamente sobre el plano bidimensional.

El archivo binario \texttt{.pose} conserva las tres fases fonológicas estándar. Para aislar el tramo con valor morfosintáctico distintivo sin alterar el archivo original, el sistema registra los marcadores temporales de inicio ($t_{\text{inicio}}$) y fin ($t_{\text{fin}}$) del núcleo en un índice de anotaciones externo. Durante la renderización, la lectura consulta este índice para aislar el segmento relevante.

\subsection{Preparación, acondicionamiento y ensamblado cinemático}
Antes de la reproducción secuencial, cada clip atraviesa una etapa de acondicionamiento estructurada:
\begin{enumerate}
    \item \textbf{Normalización temporal:} Conversión de los clips a una cadencia común de 30 fotogramas por segundo.
    \item \textbf{Aislamiento del núcleo:} Recorte del tramo delimitado por $t_{\text{inicio}}$ y $t_{\text{fin}}$.
    \item \textbf{Interpolación de faltantes:} Relleno de ausencias puntuales mediante interpolación entre fotogramas válidos adyacentes.
    \item \textbf{Ajuste de extremidad inactiva:} Reposicionamiento de la mano no participante hacia el costado del torso.
    \item \textbf{Filtro de suavizado:} Aplicación de un promedio exponencial para atenuar la oscilación entre fotogramas.
    \item \textbf{Ajuste de escala:} Normalización del tamaño corporal en función de una distancia inter-hombros fija.
\end{enumerate}

Para evitar discontinuidades espaciales abruptas al empalmar el núcleo del gesto $k$ con el núcleo $k+1$, el sistema interpola un \textit{puente sintético} de fotogramas de transición mediante la función de suavizado \textit{smoothstep}:
\begin{equation}
S(x) = 3x^2 - 2x^3 \quad \text{para } x \in [0, 1]
\label{eq:smoothstep}
\end{equation}
Este puente ejecuta una traslación suave sobre las coordenadas planares. Se explicita que este procedimiento constituye un ajuste gráfico e interpolatorio, diferendo del fenómeno biomecánico natural de la coarticulación gestual. El Cuadro~\ref{tab:salvaguardas_pose} sintetiza las previsiones de diseño ante contingencias en la representación.

\begin{table}[H]
  \centering
  \small
  \caption{Previsiones del diseño frente a fallos de estimación o léxico.}
  \label{tab:salvaguardas_pose}
  \begin{tabularx}{\textwidth}{@{}>{\raggedright\arraybackslash}p{0.42\textwidth} >{\raggedright\arraybackslash}X@{}}
    \toprule
    \textbf{Situación} & \textbf{Previsión} \\
    \midrule
    Puntos nulos puntuales en el estimador. & Interpolación lineal a partir de fotogramas válidos vecinos. \\
    \addlinespace[2pt]
    Inexistencia de archivo para una glosa. & Activación del subsistema de deletreo manual dactilológico. \\
    \addlinespace[2pt]
    Ausencia de archivos de deletreo. & Mantenimiento del esqueleto en postura estática de reposo e inclusión de aviso en interfaz. \\
    \addlinespace[2pt]
    Secuencias extensas ($>24$ glosas o $>1200$ cuadros). & Truncado de la secuencia en el límite máximo tolerable. \\
    \bottomrule
  \end{tabularx}
\end{table}

\subsection{Estrategia de contingencia léxica}
Ante palabras ausentes en el vocabulario grabado (nombres propios, neologismos o términos técnicos), el sistema aplica un mecanismo de descomposición dactilológica. La palabra se segmenta en sus letras constituyentes y se sintetiza mediante la secuencia correspondiente de caracteres del alfabeto dactilológico, consultando el endpoint específico \texttt{/fingerspell}.

\section{Diseño de la extensión de videollamada}
\label{sec:diseno_extension}

La integración con la plataforma Google Meet se resuelve mediante una extensión de navegador que intercepta la señal de video y gestiona la superposición de capas gráficas. La Figura~\ref{fig:arquitectura_video_ilsa} detalla la distribución funcional de sus componentes.

\begin{figure}[!ht]
\centering
\resizebox{1\textwidth}{!}{%
\begin{circuitikz}
\tikzstyle{every node}=[font=\fontsize{8.0pt}{10.4pt}\selectfont]

% Marco principal
\draw [line width=0.7pt] (-3.375,15.125) rectangle (9.125,6.875);

% Cámara web
\draw (-1.875,16.75) rectangle (1.375,16);
\node[anchor=center, align=center, inner xsep=0.080cm, inner ysep=0.085cm, rounded corners=0.020cm, font=\fontsize{10.8pt}{14.1pt}\selectfont] at (-0.25,16.375) {Cámara web};

% Google Meet principal
\draw [line width=0.7pt] (-3.125,14.875) rectangle (0.875,7.125);
\node[anchor=north west, align=left, inner xsep=0.080cm, inner ysep=0.085cm, rounded corners=0.020cm, font=\fontsize{10.8pt}{14.1pt}\selectfont] at (-2.975,14.725) {Google Meet};

% HUD
\draw (2.625,7.875) rectangle (4.75,7.25);
\node[anchor=center, align=center, inner xsep=0.080cm, inner ysep=0.085cm, rounded corners=0.020cm, font=\fontsize{10.8pt}{14.1pt}\selectfont] at (3.6875,7.5625) {HUD};

% Offscreen
\draw [line width=0.7pt] (2.5,14.75) rectangle (7,10.625);
\node[anchor=north east, align=right, inner xsep=0.080cm, inner ysep=0.085cm, rounded corners=0.020cm, font=\fontsize{10.8pt}{14.1pt}\selectfont] at (6.85,14.6) {offscreen};

% ILSA (motor local)
\draw (-3.25,4.75) rectangle (9.125,4);
\node[anchor=center, align=center, inner xsep=0.080cm, inner ysep=0.085cm, rounded corners=0.020cm, font=\fontsize{10.8pt}{14.1pt}\selectfont] at (2.9375,4.375) {ILSA (motor local)};

% Cajas dentro de Google Meet
\draw (-2.875,12.375) rectangle (0.125,11.25);
\node[anchor=center, align=center, inner xsep=0.080cm, inner ysep=0.085cm, rounded corners=0.020cm, font=\fontsize{10.8pt}{14.1pt}\selectfont] at (-1.375,11.8125) {video oculto \\ (flujo real)};

\draw (-2.875,14) rectangle (0.125,12.875);
\node[anchor=center, align=center, inner xsep=0.080cm, inner ysep=0.085cm, rounded corners=0.020cm, font=\fontsize{10.8pt}{14.1pt}\selectfont] at (-1.375,13.4375) {inject-gum.js \\ (toma la camara)};

\draw (-2.875,10.25) rectangle (0.625,9.125);
\node[anchor=center, align=center, inner xsep=0.080cm, inner ysep=0.085cm, rounded corners=0.020cm, font=\fontsize{10.8pt}{14.1pt}\selectfont] at (-1.125,9.6875) {canvas virtual \\ (imagen y subtítulo)};

\draw (-2.875,8.5) rectangle (0.125,7.375);
\node[anchor=center, align=center, inner xsep=0.080cm, inner ysep=0.085cm, rounded corners=0.020cm, font=\fontsize{10.8pt}{14.1pt}\selectfont] at (-1.375,7.9375) {Google Meet \\ (videollamada)};

% Cajas dentro de Offscreen
\draw (2.75,14) rectangle (6.75,12.75);
\node[anchor=center, align=center, inner xsep=0.080cm, inner ysep=0.085cm, rounded corners=0.020cm, font=\fontsize{10.8pt}{14.1pt}\selectfont] at (4.75,13.375) {controlador \\ (fotogramas reducidos)};

\draw (2.75,11.75) rectangle (6.75,10.875);
\node[anchor=center, align=center, inner xsep=0.080cm, inner ysep=0.085cm, rounded corners=0.020cm, font=\fontsize{10.8pt}{14.1pt}\selectfont] at (4.75,11.3125) {temporizador};

% Micrófono o chat
\draw (4.625,6.625) rectangle (9.125,5.875);
\node[anchor=center, align=center, inner xsep=0.080cm, inner ysep=0.085cm, rounded corners=0.020cm, font=\fontsize{10.8pt}{14.1pt}\selectfont] at (6.875,6.25) {Micrófono o chat};

% Sandbox
\draw [line width=0.7pt] (4.25,10.25) rectangle (8.75,8.375);
\node[anchor=north west, align=left, inner xsep=0.080cm, inner ysep=0.085cm, rounded corners=0.020cm, font=\fontsize{10.8pt}{14.1pt}\selectfont] at (4.4,10.1) {sandbox};

\draw (4.5,9.5) rectangle (8.5,8.625);
\node[anchor=center, align=center, inner xsep=0.080cm, inner ysep=0.085cm, rounded corners=0.020cm, font=\fontsize{10.8pt}{14.1pt}\selectfont] at (6.5,9.0625) {extracción de puntos};

% Conexiones y flechas
\draw [-{Stealth[scale=1.5]}] (0,16) -- (0,14);
\draw [-{Stealth[scale=1.5]}] (-1.25,12.875) -- (-1.25,12.375);
\draw [-{Stealth[scale=1.5]}] (-1.375,11.25) -- (-1.375,10.25) node[pos=0.5, fill=white, fill opacity=1, text opacity=1, inner xsep=0.080cm, inner ysep=0.085cm, rounded corners=0.020cm]{frame x frame};
\draw [-{Stealth[scale=1.5]}] (-1.375,9.125) -- (-1.375,8.5);

\draw [-{Stealth[scale=1.5]}] (0.125,11.875) -- (1.625,11.875) -- (1.625,13.5) -- (2.75,13.5);
\node[fill=white, fill opacity=1, text opacity=1, inner xsep=0.080cm, inner ysep=0.085cm, rounded corners=0.020cm, font=\fontsize{8.0pt}{10.4pt}\selectfont] at (1.625,12.5) {JPEG};

\draw [-{Stealth[scale=1.5]}] (6.75,13.875) -- (8.25,13.875) -- (8.25,9.5);
\node[fill=white, fill opacity=1, text opacity=1, inner xsep=0.080cm, inner ysep=0.085cm, rounded corners=0.020cm, font=\fontsize{8.0pt}{10.4pt}\selectfont] at (8.25,12.5) {frame};

\draw [dash pattern=on 6pt off 2pt, -{Stealth[scale=1.5]}] (7.625,9.5) -- (7.625,13.25) -- (6.75,13.25);
\node[fill=white, fill opacity=1, text opacity=1, inner xsep=0.080cm, inner ysep=0.085cm, rounded corners=0.020cm, font=\fontsize{8.0pt}{10.4pt}\selectfont] at (7.625,11.875) {puntos};

\draw [-{Stealth[scale=1.5]}] (2.75,11.25) -- (1.625,11.25) -- (1.625,4.75);
\node[fill=white, fill opacity=1, text opacity=1, inner xsep=0.080cm, inner ysep=0.085cm, rounded corners=0.020cm, font=\fontsize{8.0pt}{10.4pt}\selectfont] at (1.625,6.5) {POST /sign};

\draw [-{Stealth[scale=1.5]}] (5,5.875) -- (5,4.75);
\node[fill=white, fill opacity=1, text opacity=1, inner xsep=0.080cm, inner ysep=0.085cm, rounded corners=0.020cm, font=\fontsize{8.0pt}{10.4pt}\selectfont] at (5,5.375) {POST /stt};

\draw [dash pattern=on 6pt off 2pt, -{Stealth[scale=1.5]}] (6.5,4.75) -- (6.5,5.875);
\node[fill=white, fill opacity=1, text opacity=1, inner xsep=0.080cm, inner ysep=0.085cm, rounded corners=0.020cm, font=\fontsize{8.0pt}{10.4pt}\selectfont] at (6.5,5.375) {oración};

\draw [-{Stealth[scale=1.5]}] (8.375,5.875) -- (8.375,4.75);
\node[fill=white, fill opacity=1, text opacity=1, inner xsep=0.080cm, inner ysep=0.085cm, rounded corners=0.020cm, font=\fontsize{8.0pt}{10.4pt}\selectfont] at (8.375,5.375) {POST /hearing};

\draw [-{Stealth[scale=1.5]}] (3.375,4.75) -- (3.375,7.25);
\node[fill=white, fill opacity=1, text opacity=1, inner xsep=0.080cm, inner ysep=0.085cm, rounded corners=0.020cm, font=\fontsize{8.0pt}{10.4pt}\selectfont] at (3.375,6) {estado / glosas};

\draw [dash pattern=on 6pt off 2pt, -{Stealth[scale=1.5]}] (-2.875,4.75) -- (-2.875,5.75) -- (-4.75,5.75) -- (-4.75,9.75) -- (-2.875,9.75);
\node[fill=white, fill opacity=1, text opacity=1, inner xsep=0.080cm, inner ysep=0.085cm, rounded corners=0.020cm, font=\fontsize{8.0pt}{10.4pt}\selectfont] at (-3.875,5.75) {animación};

\draw [dash pattern=on 6pt off 2pt, -{Stealth[scale=1.5]}] (-2.375,4.75) -- (-2.375,6.375) -- (-4.125,6.375) -- (-4.125,9.375) -- (-2.875,9.375);
\node[fill=white, fill opacity=1, text opacity=1, inner xsep=0.080cm, inner ysep=0.085cm, rounded corners=0.020cm, font=\fontsize{8.0pt}{10.4pt}\selectfont] at (-3.25,6.375) {subtítulos};

\draw [-{Stealth[scale=1.5]}] (3.625,12.75) -- (3.625,11.75);
\node[fill=white, fill opacity=1, text opacity=1, inner xsep=0.080cm, inner ysep=0.085cm, rounded corners=0.020cm, font=\fontsize{8.0pt}{10.4pt}\selectfont] at (3.625,12.375) {puntos};

\end{circuitikz}%
}
\caption{Diagrama de arquitectura del flujo de video, procesamiento offscreen e integración con el motor ILSA.}
\label{fig:arquitectura_video_ilsa}
\end{figure}

\section{Patrones de diseño de software aplicados}
El diseño incorpora patrones de software consolidados para organizar las responsabilidades del sistema:
\begin{itemize}
    \item \textbf{Estrategia (\textit{Strategy}):} Aplicado en el módulo de filtrado de repeticiones (\texttt{RepeatGate}). Define políticas diferenciadas según la clase léxica (dígitos admiten repetición continua, letras aceptan hasta dos instancias consecutivas y vocabulario general descarta duplicados inmediatos).
    \item \textbf{Fachada (\textit{Facade}):} Aplicado en la API HTTP del motor ILSA local. Expone una interfaz unificada que oculta la complejidad interna de extracción, clasificación y comunicación con la nube.
    \item \textbf{Adaptador (\textit{Adapter}):} Implementado en la extensión web y en los conectores de red para transformar los formatos nativos del navegador a los contratos JSON requeridos por los servicios REST.
    \item \textbf{Estado (\textit{State}):} Gestiona la modalidad operativa de la sesión (Señante u Oyente), alterando dinámicamente el comportamiento de los endpoints sin requerir instancias de servicio separadas.
\end{itemize}

\section{Diseño conceptual de la arquitectura escalable de producción en la nube}
\label{sec:arq_escalable}

Para un escenario de producción con múltiples llamadas concurrentes, el diseño proyecta un esquema con balanceador de carga, réplicas sin estado de los traductores semánticos y distribución de archivos mediante almacenamiento de objetos con red de entrega de contenidos (CDN).